problem:  
Let \((M^n,g)\) be a complete, noncompact Riemannian manifold with \(\mathrm{Ric}_M \ge 0\) and **unique asymptotic cone** \((C_\infty,o)\) in the pointed Gromov–Hausdorff sense after scaling by \(R^{-2}g\) as \(R\to\infty\).  
Let \((N^m,h)\) be a compact Riemannian manifold with strictly negative sectional curvature, and \(f_0:(M,g)\to(N,h)\) a finite‑energy smooth map.  

Let \(f_t\) be the global solution to the harmonic map heat flow
\[
\partial_t f_t = \tau(f_t), \qquad f_0 = f.
\]
Assume \(E(f_t)=\int_M|df_t|^2\,d\mathrm{vol}_g\) stays finite for all \(t>0\).  
For each scaling factor \(R>0\), let \(g_R=g/R^2\) and consider the time \(t=R^2.\)  Passing to a sequence \(R_i\to\infty\), we obtain convergence \((M,g_{R_i},p)\to (C_\infty,o)\) and, after weak Sobolev limits in the Korevaar–Schoen sense, a limit map
\[
f_\infty^\ast:(C_\infty,o)\longrightarrow (N,h).
\]
We call \(f_\infty^\ast\) a **tangent harmonic map at infinity** of the heat flow.

> **Question:**  
> (1) Under these assumptions, is \(f_\infty^\ast\) unique—independent of all subsequence choices \(R_i\)—provided the asymptotic cone of \(M\) is unique?  
>  
> (2) If unique, does \(f_\infty^\ast\) minimize the Korevaar–Schoen energy in its trace class at infinity on \(C_\infty\)?  
>  
> Equivalently, does the harmonic map heat flow on a \(\mathrm{Ric}\ge0\) manifold with a unique cone at infinity canonically determine a single asymptotic harmonic map on that cone, independent of initial data modulo decay at infinity?

Why is it a "good" problem:  
This sharpened version removes the earlier ambiguities and rests purely on well‑developed, canonical frameworks—Ricci‑limit geometry, harmonic map heat flow, and metric‑space Sobolev analysis—while touching a conceptually deep unknown.

1. **Natural intersection of core theories.**  
   The problem synchronizes two fundamental geometric structures:  
   - Parabolic evolution of maps via the harmonic map heat flow, and  
   - Asymptotic geometric limits of nonnegatively Ricci curved spaces (Cheeger–Colding theory).  
   Whether the flow induces a *unique* tangent harmonic map on the asymptotic cone probes the continuity of analytic behavior under geometric scaling limits.

2. **Borderline rigidity and canonicity.**  
   Uniqueness of blow‑downs for Einstein or minimal structures under Ric ≥ 0 is extremely subtle; even tangent cones themselves can be nonunique.  Here, assuming a unique cone, the question asks whether the **analytic flow** inherits that uniqueness on the level of harmonic maps. A positive result would reveal a new parabolic rigidity principle connecting curvature, flow, and metric limits.

3. **Analytic–geometric bridge through Sobolev limit theory.**  
   The Korevaar–Schoen framework allows one to define energy‑minimizing maps into NPC targets on possibly singular spaces like cones. Thus the problem connects continuous geometric analysis on smooth manifolds to weakly convergent analysis on singular Ricci‑limit cones, demanding truly cross‑disciplinary methods.

4. **Simplicity and profound depth.**  
   The statement can be phrased in one line—“Does harmonic map heat flow on a Ric ≥ 0 manifold with a unique asymptotic cone have a unique tangent harmonic limit on that cone?”—yet solving it would require profound new insights about long‑time heat‑flow limits, energy concentration, and metric‑space functionals.  

In short, the problem is logically coherent, naturally grounded in existing theory, and exposes an uncharted connection between parabolic dynamics and global Ricci‑limit geometry.  It has the essential hallmarks of a “good” research problem: clear formulation, deep structural significance, and cross‑field reach.